\chapter*{How To Use This Book}
\markboth{How To Use This Book}{How To Use This Book}
\phantomsection
\addcontentsline{toc}{chapter}{How To Use This Book}

You can use this book in two ways. You can mount \diskname\ and run the
programs that are already on the disk, or you can type the listings yourself.
Typing is slower, but it teaches you more.

\section*{If You Want To Run A Program}
\phantomsection
\addcontentsline{toc}{section}{If You Want To Run A Program}

Mount \diskname, start NovaBASIC, and type:

\begin{lstlisting}
DIR
LOAD "TINY-1"
RUN
\end{lstlisting}

Use \texttt{DIR} when you want to see the programs on the disk. Use
\texttt{LIST} when you want to see the lines inside a program.

\section*{If You Want To Type A Program}
\phantomsection
\addcontentsline{toc}{section}{If You Want To Type A Program}

Start with \commonlisting. It begins at line 60000 and gives every game a
title screen. Type it once, save it, and keep it.

\begin{lstlisting}
NEW
REM TYPE COMMON.BAS HERE
SAVE "COMMON"
\end{lstlisting}

Then type a game listing. Game listings begin near line 10 and call the common
framework with \texttt{GOSUB 60000}.

\begin{typeinbox}
NovaBASIC only cares about the first two characters of a variable name. This
book uses short names on purpose. \texttt{SC} is a good score variable.
\texttt{SCORE} looks nicer, but NovaBASIC treats it like \texttt{SC}.
\end{typeinbox}

\section*{How To Check Your Typing}
\phantomsection
\addcontentsline{toc}{section}{How To Check Your Typing}

When a program will not run, do not erase it and start over. Check it in small
pieces.

\begin{enumerate}
\item Use \texttt{LIST} to look at the lines near the error.
\item Check every quote mark in a \texttt{PRINT} line.
\item Check every comma, colon, and parenthesis.
\item Make sure \texttt{FOR} has a matching \texttt{NEXT}.
\item Make sure \texttt{GOSUB} has a matching \texttt{RETURN}.
\end{enumerate}

\section*{Make It Yours}
\phantomsection
\addcontentsline{toc}{section}{Make It Yours}

After a program works, change it. Try a new color, a different sound, a faster
enemy, or a bigger score. You learn the machine by giving it your own ideas.
